\honest{Rationality: An Introduction}{Rob Bensinger}{2015}

In 1951 a rough football game between Dartmouth and Princeton was studied by two psychologists. ``86\% of Princeton students believed that Dartmouth had started it, whereas only 36\% of Dartmouth students blamed Dartmouth.'' \nb{Another 53\% of Dartmouth students said both teams started it, as the introduction mentions in a parenthesis. So 89\% of them blamed their own team at least in part, against 97\% at Princeton.} Shown a film of the game, Princeton students counted a mean of 9.8 Dartmouth infractions and Dartmouth students 4.3. When something we value is threatened, our perceptions and our reasoning rally to its defense; some psychologists think our ability to give reasons evolved specifically to win arguments. The stories we tell about our own reasons are often confabulated, and we trust stories for their appeal rather than for how well they predict.

Where can we find firmer ground? In mathematics. Axioms for arithmetic let us say why 2 + 2 = 4, not just that it seems right; in the same way probability theory fixes the best beliefs given priors and evidence, and decision theory the best action given beliefs and preferences. Suppose one of six classmates is your secret admirer; with no reason to favour anyone, the odds that it is Bob are 1:5, not fifty-fifty. If people with a crush wink at you ten times as often, and Bob winks, it would be a mistake to stay skeptical and a mistake to call it certain. I multiply 1:5 by 10:1 and get 2:1, a probability of two thirds; any other confidence would be inconsistent. \nb{The arithmetic is right. One step is loosely described: the introduction calls the winking ratio ``the 10:1 odds in favor of `a random person who winks at me has a crush on me'\,''. It is a likelihood ratio, as the introduction says two paragraphs later. The odds that a random winker has a crush also depend on how many people have one.}

So ``the question `What should I believe?' has an objectively right answer,'' and ``There is always a correct amount of confidence to have in a statement.'' \nb{The example has a right answer because it supplies its own starting odds and its own ratio. Consistency alone does not fix a prior, and Bayesians disagree about whether anything else does. The same point is made in the notes on ``How Much Evidence Does It Take?''} We talk as if uncertainty made belief a matter of taste. I quote Robin Hanson: ``You are never entitled to your opinion'': beliefs are about the world, not about you, there is always a best estimate, and you are entitled only to your best effort to find it. The answer about Bob is ``just as logically constrained'' as an answer in algebra.

We will never be perfect Bayesians, but the ideal shows us why an answer is right and where we went wrong. Learning arithmetic by rote, ``10 + 3 = 13'' and so on, does little good without the pattern behind it; likewise it is hard to improve your reasoning without a framework for judging whether a method works. This book is meant to help you build that framework.

The book's essays were written by Eliezer Yudkowsky for the blog \textsc{Overcoming Bias}. Scott Alexander wrote that it is useful to have as much evidence as possible, as with money, but also to use a limited amount wisely; rationality techniques get more out of the evidence we have. The sequences begin with questions where the odds are extreme and errors should be easy to spot (``Overly Convenient Excuses''); then politics, and why we reason worse, not better, about issues we call important; then rationalization, story-telling that makes beliefs feel justified without making them more accurate; then seeing evidence that does not fit our expectations; then the hazards of groups gathered around exciting ideas (``Death Spirals''); and last, ``Letting Go.''

I grant, with Luke Muehlhauser, that we ``just are cognitive biases''; biases are the ``very substance'' of our reasoning. But debiasing is not impossible: we are not perfect calculators either, yet we can train our arithmetic, learn when to trust our mathematical intuitions, and arrange our surroundings to make things easier. ``And if we're wrong today, we can be less so tomorrow.'' \nb{As in ``Biases: An Introduction,'' the claim that debiasing is possible rests on an analogy, with no evidence that teaching of this kind achieves it.}
